% ==============================================================================
% 09_criteria6_studentsupport.tex — Criterion 6: Student Support Service
% Score: 2
% ==============================================================================

\criterion{6}{บริการสนับสนุนผู้เรียน}{Student Support Service}

% ------------------------------------------------------------------------------
\criterionsub{6.1 The programme to show that the student intake policy, criteria, and procedures are clearly defined, communicated, publicly available, and up-to-date.}

มหาวิทยาลัยราชภัฏอุตรดิตถ์กำหนดนโยบาย เกณฑ์ และขั้นตอนการรับนักศึกษาอย่างชัดเจนผ่าน \textbf{ประกาศนโยบายและแนวทางการรับนักศึกษา พ.ศ.~2567} ซึ่งผ่านการเห็นชอบจากคณะกรรมการบริหารวิชาการ (กบ.วช.) สภาวิชาการ และสภามหาวิทยาลัย (อ้างอิง~\hyperlink{ref:AUNQA-6-1-1}{AUNQA-6-1-1}) โดยระบุจำนวนรับนักศึกษา เกณฑ์การรับเข้า และขั้นตอนการสมัครอย่างครบถ้วน

เกณฑ์การรับเข้าศึกษาหลักสูตร วศ.ม. วิศวกรรมคอมพิวเตอร์และปัญญาประดิษฐ์ ได้แก่ สำเร็จการศึกษาระดับปริญญาตรีสาขาที่เกี่ยวข้อง และผ่านการสัมภาษณ์โดยคณะกรรมการที่หลักสูตรแต่งตั้ง ขั้นตอนการสมัครดำเนินการผ่านระบบออนไลน์ SmartURU (\url{https://academic.uru.ac.th/SmartUru/}) ซึ่งรองรับการสมัครและติดตามสถานะออนไลน์ได้ตลอด 24 ชั่วโมง

หลักสูตรเผยแพร่ข้อมูลการรับสมัครผ่านช่องทางหลายรูปแบบ ได้แก่ (1)~เว็บไซต์มหาวิทยาลัยและคณะ (2)~Facebook Page ``วิศวกรรมคอมพิวเตอร์และปัญญาประดิษฐ์'' (495 followers ณ พฤษภาคม~2569) (3)~การประชาสัมพันธ์ตรงที่วิทยาลัยเทคนิคอุตรดิตถ์ (2 ตุลาคม~2568) และ (4)~ศูนย์ส่งเสริมบัณฑิตศึกษา มรอ. (cross-share ตุลาคม~2568)

ผลการรับนักศึกษารุ่นที่~1 (ภาค~2/2568) มีจำนวน~4 คน ซึ่งผ่านกระบวนการสัมภาษณ์และมีคุณสมบัติตามเกณฑ์ที่กำหนดครบถ้วน สำหรับการรับปีการศึกษา~2569 ภาคที่~1 มีผู้สมัครเข้ารับการสัมภาษณ์จำนวน~8 คน (ณ พฤษภาคม~2569) สะท้อนแนวโน้มการเติบโตของความต้องการอย่างต่อเนื่อง

\begin{table}[H]
  \centering
  \caption{ประวัติการรับนักศึกษา วศ.ม. CPE\&AI}
  \label{tbl:intake_history}
  \begingroup
  \footnotesize
  \setlength{\tabcolsep}{4pt}
  \renewcommand{\arraystretch}{1.2}
  \begin{tabularx}{\linewidth}{|l|c|c|l|}
    \hline
    \textbf{ปีการศึกษา / ภาค} & \textbf{ผู้สมัคร} & \textbf{รับเข้า} & \textbf{หมายเหตุ} \\
    \hline
    2568 / ภาค 2 (รุ่นแรก) & -- & 4 & เปิดรับรุ่นแรก \\
    \hline
    2569 / ภาค 1 (รุ่นสอง) & 8+ & -- (อยู่ระหว่างดำเนินการ) & สัมภาษณ์แล้ว \\
    \hline
  \end{tabularx}
  \endgroup
\end{table}

% ------------------------------------------------------------------------------
\criterionsub{6.2 The programme to show that there are adequate short-term and long-term plans for student support services.}

บริการสนับสนุนผู้เรียนครอบคลุมทั้งมิติวิชาการ (Academic Support) และมิติอื่นๆ (Non-academic Support) โดยอิงจากแผนงานระยะสั้นและระยะยาวที่ผ่านกระบวนการอนุมัติจากมหาวิทยาลัยแล้ว ดังต่อไปนี้

\begin{table}[H]
  \centering
  \caption{แผนบริการสนับสนุนผู้เรียนระยะสั้นและระยะยาว}
  \label{tbl:support_plans}
  \begingroup
  \footnotesize
  \setlength{\tabcolsep}{3pt}
  \renewcommand{\arraystretch}{1.2}
  \begin{tabularx}{\linewidth}{|>{\raggedright\arraybackslash}X|l|l|}
    \hline
    \textbf{แผน} & \textbf{ระยะ} & \textbf{หน่วยงานรับผิดชอบ} \\
    \hline
    แผนกิจกรรมพัฒนานักศึกษา & 5 ปี (2565--2569) & กองพัฒนานักศึกษา (กพน.) \\
    \hline
    แผนพัฒนาทุนวิจัยและทุนสนับสนุนการทำวิทยานิพนธ์/สารนิพนธ์ (FF, สวก., CBR, ทุนภายใน) & 5 ปี (2565--2569) & สถาบันวิจัยและพัฒนา (สวพ.) + บัณฑิตวิทยาลัย \\
    \hline
    แผนพัฒนาสมรรถนะดิจิทัล & 5 ปี (2563--2567) & สำนักวิทยบริการและเทคโนโลยีสารสนเทศ (สวท.) \\
    \hline
    แผนการจัดหา IT & ปีงบประมาณ 2568 (short-term) & สวท. \\
    \hline
  \end{tabularx}
  \endgroup
\end{table}

นอกจากแผนระดับมหาวิทยาลัยแล้ว หลักสูตรได้จัดปฐมนิเทศนักศึกษาเมื่อวันที่~12 ตุลาคม~2568 โดยมีรองอธิการบดี รศ.ดร.อิสระ อินจันทร์ และคณบดีคณะเทคโนโลยีอุตสาหกรรม ผศ.ดร.ชัชพล เกษวิริยะกิจ เข้าร่วม ซึ่งเป็นส่วนหนึ่งของแผนระยะสั้นในการให้ข้อมูลหลักสูตร นโยบายมหาวิทยาลัย และแนวทางการศึกษาระดับบัณฑิตศึกษาแก่นักศึกษาใหม่ (อ้างอิง~\hyperlink{ref:AUNQA-6-2-2}{AUNQA-6-2-2})

% ------------------------------------------------------------------------------
\criterionsub{6.3 The programme to show that there is a system to monitor, respond and report on student progress, academic performance, and workload.}

ระบบสารสนเทศสามระดับของมหาวิทยาลัยทำหน้าที่ติดตามความก้าวหน้า ผลการเรียน และภาระการเรียนของนักศึกษาแต่ละคนอย่างเชื่อมโยงกัน ดังนี้

ระบบแรกคือ ระบบอาจารย์ที่ปรึกษา (\url{https://advisor.uru.ac.th/}) (อ้างอิง~\hyperlink{ref:AUNQA-6-3-1}{AUNQA-6-3-1}) ที่ช่วยให้อาจารย์ที่ปรึกษาดูข้อมูลการลงทะเบียน ผลการเรียน และประวัติการให้คำปรึกษาของนักศึกษาแต่ละคนพร้อมบันทึกการพบปะในระบบ ระบบที่สองคือ ระบบสารสนเทศประธานหลักสูตร (\url{https://academic.uru.ac.th/teacher_chief/src/login.php}) (อ้างอิง~\hyperlink{ref:AUNQA-6-3-2}{AUNQA-6-3-2}) ที่ให้ภาพรวมผลการเรียนและสถานะการทำวิทยานิพนธ์/สารนิพนธ์ของนักศึกษาทุกคน และระบบที่สามคือ ระบบบริการนักศึกษา (\url{https://academic.uru.ac.th/std_service1/}) (อ้างอิง~\hyperlink{ref:AUNQA-6-3-3}{AUNQA-6-3-3}) ที่นักศึกษาเข้าถึงข้อมูลการลงทะเบียน ใบแจ้งหนี้ และผลการเรียนได้ด้วยตนเอง

\textbf{การติดตามวิทยานิพนธ์/สารนิพนธ์} (สำหรับหลักสูตรบัณฑิตศึกษา): อาจารย์ที่ปรึกษานัดพบนักศึกษาอย่างน้อย~1 ครั้งต่อเดือน เพื่อทบทวนความก้าวหน้าและให้ Feedback บันทึกการนัดพบผ่านระบบอาจารย์ที่ปรึกษา นักศึกษาปีการศึกษา~2568 ทั้ง~4 คนมีหัวข้อวิทยานิพนธ์/สารนิพนธ์เบื้องต้นที่อยู่ระหว่างพัฒนาให้สมบูรณ์ก่อน Proposal Defense ในวิชา 7015901/7015904 ภาคการศึกษาถัดไป ดังนี้

\begin{table}[H]
  \centering
  \caption{ความก้าวหน้าการทำวิทยานิพนธ์/สารนิพนธ์ นักศึกษารุ่นแรก}
  \label{tbl:thesis_progress}
  \begingroup
  \footnotesize
  \setlength{\tabcolsep}{3pt}
  \renewcommand{\arraystretch}{1.2}
  \begin{tabularx}{\linewidth}{|>{\raggedright\arraybackslash}p{3cm}|>{\raggedright\arraybackslash}X|l|}
    \hline
    \textbf{นักศึกษา} & \textbf{หัวข้อวิจัย} & \textbf{สาขาหลัก} \\
    \hline
    นายเกริกเกียรติ บุญไทย & พัฒนาโมเดลตรวจจับผู้ขับขี่ไม่สวมหมวกนิรภัยด้วย Deep Learning & Computer Vision \\
    \hline
    นางสาวพอใจ สุขหา & ตู้อบแห้งข้าวแต๋นระบบไฮบริดอัจฉริยะด้วย IoT เพื่อผลิตภัณฑ์ชุมชน & IoT / Smart System \\
    \hline
    นายจาตุรงค์ ทรงภักดี & การพยากรณ์หน่วยสูญเสียในสถานี EV Charging ด้วย AI ควบคุม Capacitor Bank (PEA) & AI / Power System \\
    \hline
    นายจิรกิตติ์ วราภรณ์ & การระบุพื้นที่ปลูกอ้อยด้วยภาพถ่ายดาวเทียม Sentinel-2 โดยใช้ Deep Learning & Remote Sensing \\
    \hline
  \end{tabularx}
  \endgroup
\end{table}

เมื่อนักศึกษาติดปัญหาหรือความก้าวหน้าช้ากว่าแผน อาจารย์ที่ปรึกษาให้ Feedback และดำเนินการ Corrective Actions โดยให้คำแนะนำเพิ่มเติม และกรณีจำเป็นจะส่งเรื่องให้คณะกรรมการบริหารหลักสูตรพิจารณาให้การสนับสนุนพิเศษ เช่น การปรับแผนวิจัย หรือแนะนำแหล่งข้อมูลเพิ่มเติม

% ------------------------------------------------------------------------------
\criterionsub{6.4 The programme to show that there are co-curricular activities, student competitions, and other support services to enhance the learning experience and employability of students.}

นักศึกษาในหลักสูตรมีโอกาสเข้าร่วมกิจกรรมที่เสริมประสบการณ์การเรียนรู้และพัฒนา employability ทั้งในรูปแบบกิจกรรมเสริมหลักสูตรและการมีส่วนร่วมกับชุมชน กิจกรรมเสริมหลักสูตร (Co-curricular Activities) ในปีการศึกษา~2568 เริ่มจากการปฐมนิเทศนักศึกษาใหม่ (12 ตุลาคม~2568) โดยรองอธิการบดีและคณบดี ซึ่งให้ข้อมูลเส้นทางอาชีพและแนวทางการศึกษาระดับบัณฑิตศึกษา ตามด้วย Guest Lecture Series โดยผู้เชี่ยวชาญจากภายนอก ได้แก่ อ.ธีรศักดิ์ อุปการ บรรยายหัวข้อ ``โดรน AI ตรวจสอบหลังคา'' (12 มกราคม~2569) และ อ.พิภพ มณีจำนง มรอ.วิทยาเขตน่าน บรรยาย Special Class เทคโนโลยีตู้อบ (29 มกราคม~2569) รวมทั้ง Workshop Backward Design ร่วมกับ อ.คชรัตน์ ภูฆัง (คณะครุศาสตร์) และ ผศ.ดร.กชกร ลาภมาก (คณะวิทยาศาสตร์) (19 มกราคม~2569) เพื่อพัฒนาทักษะการออกแบบการเรียนการสอนแบบข้ามสาขา

ในด้านกิจกรรมชุมชน (Community Engagement as Learning) นักศึกษาบัณฑิตศึกษาได้รับโอกาสเป็นผู้ให้ความรู้แก่สังคม ซึ่งเสริมทักษะการสื่อสาร ความเป็นผู้นำ และการถ่ายทอดองค์ความรู้ในสถานการณ์จริง โดยทำหน้าที่วิทยากรสอน AI ให้ครูโรงเรียนเทศบาลสามแห่ง ได้แก่ โรงเรียนเทศบาลท่าอิฐ วัดหนองผา และวัดเกษมจิตตาราม (26 มีนาคม~2569) ทั้งยังมี GenAI Short Courses รวม~4 คอร์ส (เมษายน--พฤษภาคม~2569) สำหรับบุคลากรการศึกษา หลักสูตรใช้ Facebook Page ``วิศวกรรมคอมพิวเตอร์และปัญญาประดิษฐ์'' เป็นช่องทางหลักในการบันทึกและเผยแพร่กิจกรรมตลอดปีการศึกษา~2568 มีโพสต์รวม~99 รายการ ครอบคลุมกิจกรรมในชั้นเรียน งานวิชาการ และการมีส่วนร่วมกับชุมชน

ในระดับมหาวิทยาลัย นักศึกษาสามารถเข้าร่วมกิจกรรมพัฒนานักศึกษาที่จัดโดยองค์การบริหารกิจกรรมนักศึกษา สโมสรนักศึกษา และกองพัฒนานักศึกษา (กพน.) เช่น กีฬามหาวิทยาลัยแห่งประเทศไทย ครั้งที่~51 กีฬาราชภัฏภาคเหนือ ครั้งที่~35 และโครงการ TO BE NUMBER ONE ด้านผลการแข่งขัน นายจิรกิตติ์ วราภรณ์ (นักศึกษา รหัส~68 ครูผู้ช่วย วิทยาลัยเทคนิคอุตรดิตถ์) ได้รับรางวัลเหรียญทองแดงระดับภาคเหนือ (มกราคม~2569) และรางวัลชนะเลิศอันดับ~1 ในงานวันเทคโนโลยี ครั้งที่~11 (กุมภาพันธ์~2569) แสดงให้เห็นว่ากิจกรรมของหลักสูตรส่งผลให้นักศึกษามีผลงานที่ได้รับการยอมรับในระดับภาค

% ------------------------------------------------------------------------------
\criterionsub{6.5 The programme to show that the competences of support staff are defined, assessed and communicated.}

มหาวิทยาลัยราชภัฏอุตรดิตถ์กำหนดสมรรถนะและบทบาทของเจ้าหน้าที่สายสนับสนุนที่ให้บริการผู้เรียนอย่างชัดเจน โดยแบ่งตามประเภทงาน ณ วันที่ 30 เมษายน~2569 มีเจ้าหน้าที่สนับสนุนรวม \textbf{143 คน} ดังตาราง~\ref{tbl:support_staff}

\begin{table}[H]
  \centering
  \caption{จำนวนและคุณวุฒิเจ้าหน้าที่สนับสนุนผู้เรียน (ณ 30 เมษายน 2569)}
  \label{tbl:support_staff}
  \begingroup
  \footnotesize
  \setlength{\tabcolsep}{3pt}
  \renewcommand{\arraystretch}{1.2}
  \begin{tabularx}{\linewidth}{|>{\raggedright\arraybackslash}X|c|c|c|c|}
    \hline
    \textbf{ประเภทงาน} & \textbf{มัธยม} & \textbf{ป.ตรี} & \textbf{ป.โท} & \textbf{รวม} \\
    \hline
    รับสมัครนักศึกษา (กบศ.) & 2 & 7 & 8 & 17 \\
    \hline
    จัดการเรียนการสอน (กบศ./สหกิจ/วิเทศ/LLL/GE) & 2 & 17 & 14 & 33 \\
    \hline
    วิจัย (สวพ.) & 5 & 6 & 5 & 16 \\
    \hline
    บริการวิชาการ (ศวท./สวพ./อวน.) & 12 & 10 & 7 & 29 \\
    \hline
    กิจกรรมเสริมหลักสูตร (กพน./สวส.) & 9 & 28 & 11 & 48 \\
    \hline
    \textbf{รวมทั้งสิ้น} & \textbf{30} & \textbf{68} & \textbf{45} & \textbf{143} \\
    \hline
  \end{tabularx}
  \endgroup
\end{table}

เจ้าหน้าที่สนับสนุนได้รับการประเมินสมรรถนะตามกรอบที่กำหนด และผ่านการพัฒนาศักยภาพอย่างต่อเนื่องผ่านการฝึกอบรม เพื่อให้บริการผู้เรียนได้อย่างมีประสิทธิภาพ (อ้างอิง~\hyperlink{ref:AUNQA-6-5-1}{AUNQA-6-5-1} ถึง~6-5-3)

% ------------------------------------------------------------------------------
\criterionsub{6.6 The programme to show that student support services are evaluated and benchmarked for improvement.}

คุณภาพบริการสนับสนุนผู้เรียนได้รับการประเมินผ่านแบบสำรวจความพึงพอใจของนักศึกษาและการเปรียบเทียบกับสถาบันอ้างอิง โดยผลการประเมินในปีการศึกษา~2568 แบ่งออกเป็นสองด้านหลัก ได้แก่ ผลประเมินระบบรับนักศึกษา (ตาราง~\ref{tbl:admission_satisfaction}) และผลประเมินบริการสนับสนุนผู้เรียน (ตาราง~\ref{tbl:student_service_satisfaction}) ซึ่งทั้งสองด้านให้คะแนนเฉลี่ยมากกว่า~4.5 จาก~5.0 สะท้อนว่าบริการสนับสนุนผู้เรียนของมหาวิทยาลัยมีคุณภาพและตอบสนองความต้องการของผู้เรียนได้ดี
\begin{table}[H]
  \centering
  \caption{ผลประเมินความพึงพอใจการรับนักศึกษา}
  \label{tbl:admission_satisfaction}
  \begingroup
  \footnotesize
  \renewcommand{\arraystretch}{1.2}
  \begin{tabular}{|l|c|}
    \hline
    \textbf{หมวด} & \textbf{คะแนน (5 คะแนน)} \\
    \hline
    ประชาสัมพันธ์/ข้อมูล & 4.63 \\
    \hline
    กระบวนการคัดเลือก & 4.64 \\
    \hline
    ระบบ SmartURU & 4.52 \\
    \hline
    บริการเจ้าหน้าที่ & 4.54 \\
    \hline
    \textbf{ภาพรวม} & \textbf{4.58} \\
    \hline
  \end{tabular}
  \endgroup
\end{table}

\begin{table}[H]
  \centering
  \caption{ผลประเมินความพึงพอใจบริการสนับสนุนผู้เรียน}
  \label{tbl:student_service_satisfaction}
  \begingroup
  \footnotesize
  \renewcommand{\arraystretch}{1.2}
  \begin{tabular}{|l|c|}
    \hline
    \textbf{หมวด} & \textbf{คะแนน (5 คะแนน)} \\
    \hline
    ประสิทธิภาพการสอน & 4.56 \\
    \hline
    กิจกรรมพัฒนา & 4.49 \\
    \hline
    การให้คำปรึกษา & 4.56 \\
    \hline
    ข้อมูลข่าวสาร & 4.49 \\
    \hline
    กิจกรรมวิชาการ & 4.47 \\
    \hline
    ข้อร้องเรียน & 4.46 \\
    \hline
    \textbf{ภาพรวม} & \textbf{4.51} \\
    \hline
  \end{tabular}
  \endgroup
\end{table}

หลักสูตรกำหนด Benchmarking Partners ภายนอกที่เปิดสอนหลักสูตรบัณฑิตศึกษาด้านวิศวกรรมคอมพิวเตอร์และปัญญาประดิษฐ์จริง โดยพิจารณาจากความใกล้เคียงของขอบเขตวิชาการ (CPE/AI) มากกว่าการจำกัดอยู่เฉพาะเครือข่ายมหาวิทยาลัยราชภัฏ เนื่องจากปัจจุบันยังไม่มีมหาวิทยาลัยราชภัฏใดในภาคเหนือเปิดหลักสูตรบัณฑิตศึกษาในสาขานี้ การเลือก Partners ดังตารางที่~\ref{tbl:benchmarking_plan}

\begin{table}[H]
  \centering
  \caption{แผน Formal Benchmarking บริการสนับสนุนนักศึกษา ปีการศึกษา 2569}
  \label{tbl:benchmarking_plan}
  \begingroup
  \footnotesize
  \setlength{\tabcolsep}{3pt}
  \renewcommand{\arraystretch}{1.2}
  \begin{tabularx}{\linewidth}{|>{\raggedright\arraybackslash}X|>{\raggedright\arraybackslash}X|c|c|}
    \hline
    \textbf{Benchmarking Partner} & \textbf{หลักสูตร} & \textbf{KPI ที่เปรียบเทียบ} & \textbf{กำหนดการ} \\
    \hline
    มหาวิทยาลัยเชียงใหม่ (ภาควิชาวิศวกรรมคอมพิวเตอร์) & วศ.ม. วิศวกรรมคอมพิวเตอร์ & อัตราคงอยู่, Satisfaction, Research Output & ต.ค. 2569 \\
    \hline
    สจล. (KMITL) คณะวิศวกรรมศาสตร์ & วศ.ม. วิศวกรรมไฟฟ้าและคอมพิวเตอร์ + วศ.ม. AI เพื่อวิเคราะห์ธุรกิจ & Advisor System, Time-to-graduate, Satisfaction & ต.ค. 2569 \\
    \hline
  \end{tabularx}
  \endgroup
  \begin{flushleft}
    \footnotesize KPI: (1) อัตราการคงอยู่ของนักศึกษา (Retention Rate) (2) คะแนนประเมินความพึงพอใจบริการสนับสนุน (3) จำนวนกิจกรรม co-curricular ต่อปีการศึกษา \\
    กำหนดการ: รวบรวมข้อมูล มิ.ย.--ก.ย. 2569 $\rightarrow$ วิเคราะห์และรายงานผล ต.ค. 2569 $\rightarrow$ นำผลไปปรับปรุงแผนบริการ พ.ย. 2569
  \end{flushleft}
\end{table}

ข้อมูล Baseline เบื้องต้นที่หลักสูตรรวบรวมจากแหล่งสาธารณะของ Benchmarking Partners ปีการศึกษา~2568 เพื่อใช้เป็นจุดเริ่มต้นสำหรับการเปรียบเทียบในรอบถัดไป ดังแสดงในตาราง~\ref{tbl:benchmark_baseline_2568}

\begin{table}[H]
  \centering
  \caption{Baseline ข้อมูลเปรียบเทียบเบื้องต้น --- Enrollment ปีการศึกษา 2568}
  \label{tbl:benchmark_baseline_2568}
  \begingroup
  \footnotesize
  \setlength{\tabcolsep}{4pt}
  \renewcommand{\arraystretch}{1.2}
  \begin{tabularx}{\linewidth}{|>{\raggedright\arraybackslash}X|c|c|c|>{\raggedright\arraybackslash}p{3cm}|}
    \hline
    \textbf{หลักสูตร} & \textbf{ภาคปกติ} & \textbf{ภาคพิเศษ} & \textbf{รวม} & \textbf{แหล่งข้อมูล} \\
    \hline
    \textbf{มรอ.} วศ.ม. CPE\&AI (รุ่นแรก ภาค 2/2568) & 4 & --- & \textbf{4} & ระบบรับสมัคร มรอ. \\
    \hline
    \textbf{มช.} วศ.ม. วิศวกรรมคอมพิวเตอร์ (รวมทุกชั้นปี ปี 2568) & 16 & 6 & \textbf{22} & สถิตินักศึกษา มช. รหัส 68 (\hyperlink{ref:AUNQA-6-BM-CMU-2568}{AUNQA-6-BM-CMU-2568}) \\
    \hline
    \textbf{สจล. (KMITL)} วศ.ม. วิศวกรรมไฟฟ้าและคอมพิวเตอร์ (ปี~1 ปีการศึกษา~2567 \textit{เฉพาะรุ่นใหม่}) & 39 & 3 & \textbf{42} & สถิตินักศึกษา สจล. ปี 2567 (\hyperlink{ref:AUNQA-6-BM-KMITL-2567}{AUNQA-6-BM-KMITL-2567}) \\
    \hline
    \textbf{มรภ.นครปฐม (NPRU)} วศ.ม. วิศวกรรมไฟฟ้าและปัญญาประดิษฐ์ (รวมทุกชั้นปี+ตกค้าง ปี 2568) & 8 & --- & \textbf{8} & สถิตินักศึกษา มรภ.นครปฐม ปี 2568 (\hyperlink{ref:AUNQA-6-BM-NPRU-2568}{AUNQA-6-BM-NPRU-2568}) \\
    \hline
  \end{tabularx}
  \endgroup
  \begin{flushleft}
    \footnotesize ข้อสังเกต: ข้อมูลที่แสดงเป็นนักศึกษา\textbf{ปี~1 รุ่นใหม่} (Cohort เปรียบเทียบกันได้) สำหรับ มช./สจล. และ\textbf{รวมทุกชั้นปี+ตกค้าง} สำหรับ NPRU/มรอ. ซึ่งเป็นหลักสูตรขนาดเล็ก การเปรียบเทียบในรอบถัดไป (ต.ค. 2569) จะเน้น KPI เชิงสัดส่วน (Retention~\%, Satisfaction Score, Time-to-graduate) ไม่ใช่ตัวเลขรวม
  \end{flushleft}
\end{table}

จากการเปรียบเทียบขนาดหลักสูตรกับ Partners ทั้งสามมหาวิทยาลัย พบประเด็นที่ใช้เป็นข้อมูลนำเข้าในการพัฒนาบริการสนับสนุนผู้เรียน~3 ประเด็นหลัก ประเด็นแรก หลักสูตร CPE\&AI ของ มรอ. (4 คน รุ่นแรก) มีขนาดใกล้เคียงกับ \textbf{มรภ.นครปฐม} วศ.ม. วิศวกรรมไฟฟ้าและปัญญาประดิษฐ์ (รวม 8 คน) ซึ่งเป็น peer institution ที่ใกล้เคียงที่สุดทั้งด้านประเภทมหาวิทยาลัย (ราชภัฏ) ขอบเขตวิชาการ (EE/AI/CPE) และขนาดหลักสูตร ขณะที่ มช. (22 คน) และ สจล. (42 คน ปี~1) มีขนาดใหญ่กว่า~5.5--10.5 เท่า เนื่องจากเป็นมหาวิทยาลัยวิจัยระดับชั้นนำ ดังนั้น การเปรียบเทียบเชิงสัดส่วน (Retention~\%, Satisfaction, Time-to-graduate) จึงเหมาะสมกว่าตัวเลขสัมบูรณ์

ประเด็นที่สอง เป็นที่น่าสังเกตว่า มรภ.นครปฐม \textbf{ไม่มีนักศึกษารับใหม่ในปี 2568} (ปี~1 = 0 คน, มีเพียงปี~2 จากรุ่น 2567 และนักศึกษาตกค้าง) สะท้อนความท้าทายของหลักสูตรบัณฑิตศึกษาด้าน EE/AI ในมหาวิทยาลัยภูมิภาค ซึ่งเป็น sector-wide challenge ที่หลักสูตร CPE\&AI มรอ. ควรเตรียมการรับมือ โดยเฉพาะการรักษาอัตราการรับนักศึกษาให้สม่ำเสมอผ่านการประชาสัมพันธ์เชิงรุก ความร่วมมือกับสถานประกอบการในพื้นที่ และการเสนอทุนการศึกษา/ทุนวิจัย

ประเด็นที่สาม การที่ มรอ. มีผู้สมัครเข้าสัมภาษณ์ปีการศึกษา~2569 \textbf{เพิ่มขึ้นเป็น 8 คน} (จาก 6 คนสมัครและรับ 4 คนในปี 2568) สะท้อนแนวโน้มการเติบโตของความต้องการในพื้นที่ ซึ่งเป็นสัญญาณบวกที่ \textbf{ตรงข้าม}กับ NPRU ที่ขาดรุ่นใหม่ในปี 2568 หลักสูตรจะนำผลการเปรียบเทียบนี้ไปกำหนดเป้าหมายเชิงสัดส่วนสำหรับรอบ Benchmarking ต.ค. 2569 และพัฒนาบริการสนับสนุนผู้เรียนให้สอดรับกับขนาดและบริบทของหลักสูตรขนาดเล็ก

\begin{strengthbox}
หลักสูตรมีระบบรับนักศึกษาออนไลน์ที่ใช้งานได้จริง (SmartURU) ระบบติดตามผู้เรียนครบ~3 มิติ (Progress/Performance/Workload) ผ่าน advisor.uru.ac.th และระบบประธานหลักสูตร กิจกรรมเสริมหลักสูตรที่หลากหลายมีหลักฐาน (Facebook~99 โพสต์) และผลประเมินบริการสูงกว่า~4.5/5.0 ทั้งระบบรับนักศึกษาและบริการสนับสนุน นักศึกษามีผลงานที่ได้รับรางวัลระดับภาคในปีแรก
\end{strengthbox}

\begin{weaknessbox}
Formal Benchmarking มี Baseline เปรียบเทียบกับ มช., สจล., และ มรภ.นครปฐม แล้ว (ตาราง~\ref{tbl:benchmark_baseline_2568}) แต่การเปรียบเทียบเชิงลึก (Retention, Satisfaction, Time-to-graduate) ยังกำหนดดำเนินการ ต.ค. 2569 ยังไม่มีผลจริง; บันทึก Formal Monitoring Records เฉพาะรายนักศึกษา~4 คนอาจต้องพัฒนาให้เป็นระบบมากขึ้น; Co-curricular activities เฉพาะนักศึกษาบัณฑิตศึกษา CPE\&AI ยังมีน้อยเมื่อเทียบกับกิจกรรมระดับมหาวิทยาลัย; ด้านภาษาอังกฤษ มหาวิทยาลัยมีประกาศมาตรฐานความสามารถด้านภาษาอังกฤษ พ.ศ.~2569 (AUNQA-6-1-5) แล้ว แต่ยังต้องพัฒนาแผนเสริมทักษะเฉพาะกลุ่มสำหรับนักศึกษา CPE\&AI ในรอบถัดไป
\end{weaknessbox}

\begin{selfassessmentbox}{2}{หลักสูตรมีระบบรับนักศึกษา บริการสนับสนุน และการติดตามผู้เรียนที่ชัดเจนผ่านระบบดิจิทัลของมหาวิทยาลัย ผลประเมินบริการสูง กิจกรรม co-curricular มีหลักฐาน แต่ยังต้องพัฒนา Benchmarking ภายนอกและ Formal Records รายบุคคลให้ครบถ้วน}
\end{selfassessmentbox}

\begin{evidencelist}
\evidence{AUNQA-6-1-5}{ประกาศมหาวิทยาลัยราชภัฏอุตรดิตถ์ เรื่อง มาตรฐานความสามารถด้านภาษาอังกฤษสำหรับนักศึกษา พ.ศ.~2569 (ใช้กับนักศึกษาบัณฑิตศึกษาที่เข้าศึกษาตั้งแต่ปีการศึกษา~2569 เป็นต้นไป)\ \href{https://syweb.kidcbc.work/Eviden_aunqa68/\%E0\%B8\%AB\%E0\%B8\%A5\%E0\%B8\%B1\%E0\%B8\%81\%E0\%B8\%90\%E0\%B8\%B2\%E0\%B8\%99\%E0\%B8\%AA\%E0\%B8\%B2\%E0\%B8\%A2\%E0\%B8\%AA\%E0\%B8\%99\%E0\%B8\%B1\%E0\%B8\%9A\%E0\%B8\%AA\%E0\%B8\%99\%E0\%B8\%B8\%E0\%B8\%99/AUNQA-6-1-5_\%E0\%B8\%9B\%E0\%B8\%A3\%E0\%B8\%B0\%E0\%B8\%81\%E0\%B8\%B2\%E0\%B8\%A8\%E0\%B8\%A1\%E0\%B8\%B2\%E0\%B8\%95\%E0\%B8\%A3\%E0\%B8\%90\%E0\%B8\%B2\%E0\%B8\%99\%E0\%B8\%84\%E0\%B8\%A7\%E0\%B8\%B2\%E0\%B8\%A1\%E0\%B8\%AA\%E0\%B8\%B2\%E0\%B8\%A1\%E0\%B8\%B2\%E0\%B8\%A3\%E0\%B8\%96\%E0\%B8\%A0\%E0\%B8\%B2\%E0\%B8\%A9\%E0\%B8\%B2\%E0\%B8\%AD\%E0\%B8\%B1\%E0\%B8\%87\%E0\%B8\%81\%E0\%B8\%A4\%E0\%B8\%A9\%E0\%B8\%99\%E0\%B8\%B1\%E0\%B8\%81\%E0\%B8\%A8\%E0\%B8\%B6\%E0\%B8\%81\%E0\%B8\%A9\%E0\%B8\%B2_2569.pdf}{\texttt{[PDF]}}}
\evidence{AUNQA-6-1-1}{ประกาศนโยบายและแนวทางการรับนักศึกษา พ.ศ.2567\ \href{https://syweb.kidcbc.work/Eviden_aunqa68/\%E0\%B8\%A3\%E0\%B8\%B0\%E0\%B9\%80\%E0\%B8\%9A\%E0\%B8\%B5\%E0\%B8\%A2\%E0\%B8\%9A\%E0\%B8\%82\%E0\%B9\%89\%E0\%B8\%AD\%E0\%B8\%9A\%E0\%B8\%B1\%E0\%B8\%87\%E0\%B8\%84\%E0\%B8\%B1\%E0\%B8\%9A\%E0\%B9\%81\%E0\%B8\%A5\%E0\%B8\%B0\%E0\%B8\%9B\%E0\%B8\%A3\%E0\%B8\%B0\%E0\%B8\%81\%E0\%B8\%B2\%E0\%B8\%A8\%E0\%B8\%AD\%E0\%B8\%B7\%E0\%B9\%88\%E0\%B8\%99\%E0\%B9\%86/AUNQA-6-1-1_\%E0\%B8\%99\%E0\%B9\%82\%E0\%B8\%A2\%E0\%B8\%9A\%E0\%B8\%B2\%E0\%B8\%A2\%E0\%B8\%A3\%E0\%B8\%B1\%E0\%B8\%9A\%E0\%B8\%99\%E0\%B8\%B1\%E0\%B8\%81\%E0\%B8\%A8\%E0\%B8\%B6\%E0\%B8\%81\%E0\%B8\%A9\%E0\%B8\%B2_2567.pdf}{\texttt{[PDF]}}}
\evidence{AUNQA-6-1-2}{ระบบรับนักศึกษาออนไลน์ SmartURU: \url{https://academic.uru.ac.th/SmartUru/}}
\evidence{AUNQA-6-2-2}{แผน IT 2568: \url{https://lifelong.uru.ac.th/} (\hyperlink{ref:AUNQA-6-2-2}{AUNQA-6-2-2}\_แผนIT\_2568.pdf)\ \href{https://syweb.kidcbc.work/Eviden_aunqa68/\%E0\%B8\%A3\%E0\%B8\%B0\%E0\%B9\%80\%E0\%B8\%9A\%E0\%B8\%B5\%E0\%B8\%A2\%E0\%B8\%9A\%E0\%B8\%82\%E0\%B9\%89\%E0\%B8\%AD\%E0\%B8\%9A\%E0\%B8\%B1\%E0\%B8\%87\%E0\%B8\%84\%E0\%B8\%B1\%E0\%B8\%9A\%E0\%B9\%81\%E0\%B8\%A5\%E0\%B8\%B0\%E0\%B8\%9B\%E0\%B8\%A3\%E0\%B8\%B0\%E0\%B8\%81\%E0\%B8\%B2\%E0\%B8\%A8\%E0\%B8\%AD\%E0\%B8\%B7\%E0\%B9\%88\%E0\%B8\%99\%E0\%B9\%86/AUNQA-6-2-2_\%E0\%B9\%81\%E0\%B8\%9C\%E0\%B8\%99IT_2568.pdf}{\texttt{[PDF]}}}
\evidence{AUNQA-6-3-1}{ระบบอาจารย์ที่ปรึกษาออนไลน์ มรอ.: \url{https://advisor.uru.ac.th/}}
\evidence{AUNQA-6-3-2}{ระบบสารสนเทศประธานหลักสูตร มรอ.: \url{https://academic.uru.ac.th/teacher_chief/src/login.php}}
\evidence{AUNQA-6-3-3}{ระบบบริการนักศึกษาออนไลน์ มรอ.: \url{https://academic.uru.ac.th/std_service1/}}
\evidence{AUNQA-6-6-1}{Dashboard ผลการประเมินการให้บริการหน่วยงานสนับสนุน มรอ. ปีการศึกษา~2568 \href{https://script.google.com/a/macros/uru.ac.th/s/AKfycbzxhjtd8cYhLz68UNeJavc7nuaOTgTjtKv5osom6WuQFxbYC0lRbCk2u67LVttLFhpdvw/exec}{\texttt{[Dashboard]}\ (ภาพรวม~4.51/5.0)}}
\evidence{AUNQA-6-ADM-2568}{รายงานการประชุมคณะกรรมการสัมภาษณ์รับนักศึกษา 15 พ.ย. 2568 (6 สมัคร; รับ 4 คน)\ \href{https://syweb.kidcbc.work/Eviden_aunqa68/\%E0\%B8\%AB\%E0\%B8\%A5\%E0\%B8\%B1\%E0\%B8\%81\%E0\%B8\%90\%E0\%B8\%B2\%E0\%B8\%99\%E0\%B8\%81\%E0\%B8\%B2\%E0\%B8\%A3\%E0\%B8\%9B\%E0\%B8\%A3\%E0\%B8\%B0\%E0\%B8\%8A\%E0\%B8\%B8\%E0\%B8\%A1/01_\%E0\%B8\%A3\%E0\%B8\%B0\%E0\%B8\%94\%E0\%B8\%B1\%E0\%B8\%9A\%E0\%B8\%AB\%E0\%B8\%A5\%E0\%B8\%B1\%E0\%B8\%81\%E0\%B8\%AA\%E0\%B8\%B9\%E0\%B8\%95\%E0\%B8\%A3/\%E0\%B8\%81\%E0\%B8\%A1\%E0\%B8\%A8_2568-11-15_\%E0\%B8\%AA\%E0\%B8\%B1\%E0\%B8\%A1\%E0\%B8\%A0\%E0\%B8\%B2\%E0\%B8\%A9\%E0\%B8\%93\%E0\%B9\%8C\%E0\%B8\%99\%E0\%B8\%B1\%E0\%B8\%81\%E0\%B8\%A8\%E0\%B8\%B6\%E0\%B8\%81\%E0\%B8\%A9\%E0\%B8\%B2_\%E0\%B8\%A3\%E0\%B8\%B1\%E0\%B8\%9A4\%E0\%B8\%84\%E0\%B8\%99.pdf}{\texttt{[PDF]}}}
\evidence{AUNQA-6-ADM-2569}{รายงานการประชุมสัมภาษณ์รับนักศึกษา 9 พ.ค. 2569 (8 คน)\ \href{https://syweb.kidcbc.work/Eviden_aunqa68/\%E0\%B8\%AB\%E0\%B8\%A5\%E0\%B8\%B1\%E0\%B8\%81\%E0\%B8\%90\%E0\%B8\%B2\%E0\%B8\%99\%E0\%B8\%81\%E0\%B8\%B2\%E0\%B8\%A3\%E0\%B8\%9B\%E0\%B8\%A3\%E0\%B8\%B0\%E0\%B8\%8A\%E0\%B8\%B8\%E0\%B8\%A1/01_\%E0\%B8\%A3\%E0\%B8\%B0\%E0\%B8\%94\%E0\%B8\%B1\%E0\%B8\%9A\%E0\%B8\%AB\%E0\%B8\%A5\%E0\%B8\%B1\%E0\%B8\%81\%E0\%B8\%AA\%E0\%B8\%B9\%E0\%B8\%95\%E0\%B8\%A3/\%E0\%B8\%81\%E0\%B8\%A1\%E0\%B8\%A8_2569-05-09_\%E0\%B8\%AA\%E0\%B8\%B1\%E0\%B8\%A1\%E0\%B8\%A0\%E0\%B8\%B2\%E0\%B8\%A9\%E0\%B8\%93\%E0\%B9\%8C\%E0\%B8\%99\%E0\%B8\%B1\%E0\%B8\%81\%E0\%B8\%A8\%E0\%B8\%B6\%E0\%B8\%81\%E0\%B8\%A9\%E0\%B8\%B2_\%E0\%B8\%A3\%E0\%B8\%B1\%E0\%B8\%9A8\%E0\%B8\%84\%E0\%B8\%99.pdf}{\texttt{[PDF]}}}
\evidence{AUNQA-6-BM-CMU-2568}{\hypertarget{ref:AUNQA-6-BM-CMU-2568}{}สถิติข้อมูลนักศึกษาบัณฑิตศึกษา คณะวิศวกรรมศาสตร์ มหาวิทยาลัยเชียงใหม่ ปีการศึกษา 2568 (รหัส 68): สาขาวิศวกรรมคอมพิวเตอร์ ป.โท 22 คน (ภาคปกติ 16 + ภาคพิเศษ 6), ป.เอก 4 คน --- ใช้เป็น Baseline benchmarking partner\ \href{https://syweb.kidcbc.work/Eviden_aunqa68/AUNQA-6-BM-CMU-2568_CMU_CPE_Enrollment_Stats.png}{\texttt{[PNG]}}}
\evidence{AUNQA-6-BM-KMITL-2567}{\hypertarget{ref:AUNQA-6-BM-KMITL-2567}{}สถิติข้อมูลนักศึกษาแยกตามคณะ/หลักสูตร สถาบันเทคโนโลยีพระจอมเกล้าเจ้าคุณทหารลาดกระบัง (KMITL) ปีการศึกษา 2567: ป.โท วิศวกรรมไฟฟ้าและคอมพิวเตอร์ รุ่นใหม่ (ปี~1) 42 คน (รวมทุกชั้นปี+ตกค้าง+สำเร็จ 177 คน) --- ใช้เฉพาะ ปี~1 เพื่อเทียบ cohort \textit{like-for-like}\ \href{https://syweb.kidcbc.work/Eviden_aunqa68/AUNQA-6-BM-KMITL-2567_KMITL_Graduate_Statistics.xlsx}{\texttt{[XLSX]}}}
\evidence{AUNQA-6-BM-NPRU-2568}{\hypertarget{ref:AUNQA-6-BM-NPRU-2568}{}สถิติข้อมูลนักศึกษามหาวิทยาลัยราชภัฏนครปฐม (NPRU) ปีการศึกษา 2568: ป.โท วิศวกรรมไฟฟ้าและปัญญาประดิษฐ์ (วศ.ม.) 8 คน (ปี~2/2567 = 3 คน + ตกค้าง 5 คน, ไม่มีรับใหม่ปี 2568) + ป.เอก สาขาเดียวกัน 3 คน --- เป็น peer benchmark ระดับ มรภ. ที่ใกล้เคียง CPE\&AI มรอ. มากที่สุด\ \href{https://syweb.kidcbc.work/Eviden_aunqa68/AUNQA-6-BM-NPRU-2568_NPRU_Student_Data.pdf}{\texttt{[PDF]}}}
\end{evidencelist}
